% ---- pg-189 ----
间大, 至于多大, 这就与 $p,n$ 及 $1-\alpha$ 三者的大小有关.

对 $1-\alpha=0.95, n=15$ 和 $p=4$, \hyperref[eq:5-24]{式(5-24)}及\hyperref[eq:5-27]{式(5-27)}中 $\sqrt{s_{ii}/n}$ 前的系数分别为
\[
\sqrt{\frac{p(n-1)}{(n-p)}F_{p,n-p}(0.05)}=\sqrt{\frac{4\times14}{11}}\times3.36=4.14
\]
及 $t_{n-1}(0.025)=2.145$. 因此这种情况下的联合置信区间要比单一 $t$ 区间宽 $100\times(4.14-2.145)/2.145=93\%$.

\hyperref[tab:5.3]{表 5.3} 给出了按\hyperref[eq:5-21]{式(5-21)}算出的单一 $t$ 区间及相应的联合 $T^{2}$ 区间的距离乘数的临界值. 一般来说, 相对 $t$ 区间而言, $T^{2}$ 区间的宽度当 $n$ 固定而 $p$ 增大时增大; 当 $p$ 固定而 $n$ 增大时减小.

\begin{mvtable}
\phantomsection\label{tab:5.3}\mvtabcap{表 5.3\quad 对选定的 $n$ 与 $p$ ($1-\alpha=0.95$) 单一 $t$ 区间与 $T^{2}$ 区间的距离乘数之临界值}
\begin{tabular}{rccc}
\toprule
\multirow{2}{*}{$n$} & \multirow{2}{*}{$t_{n-1}(0.025)$} & \multicolumn{2}{c}{$\sqrt{\dfrac{(n-1)p}{(n-p)}F_{p,n-p}(0.05)}$}\\
\cmidrule(lr){3-4}
& & $p=4$ & $p=10$\\
\midrule
15 & 2.145 & 4.14 & 11.52\\
25 & 2.064 & 3.60 & 6.39\\
50 & 2.010 & 3.31 & 5.05\\
100 & 1.970 & 3.19 & 4.61\\
$\infty$ & 1.960 & 3.08 & 4.28\\
\bottomrule
\end{tabular}
\end{mvtable}

将\hyperref[tab:5.3]{表 5.3} 中的置信区间进行比较显然有些不公平, 因为对固定的 $n$ 与 $p$, 其所有的 $T^{2}$ 区间相联系的置信水平为 0.95, 而在同样的 $n$ 值, 所有单一 $t$ 区间的置信水平却比 0.95 要小些. 单一 $t$ 区间的宽度太短, 以至于对所有 $p$ 个均值的单独置信区间不能有一个总的置信水平. 然而, 当这是唯一可作的推断时, 我们有时就把它们看作是均值的最佳可能信息. 此外, 如果仅当 $T^{2}$ 检验拒绝零假设时才计算单一置信区间, 那么有些研究者认为, 单一置信区间也许比 $T^{2}$ 区间更精确地代表有关均值的信息.

如果 $T^{2}$ 区间仅应用于 $p$ 个均值分量, 那么 $T^{2}$ 区间显然太宽. 为了说明其原因, 我们来考虑\hyperref[fig:5.2]{图 5.2} 中的置信椭圆及联合置信区间. 若 $\mu_{1},\mu_{2}$ 分别落在各自的 $T^{2}$ 区间内, 那么 $(\mu_{1},\mu_{2})$ 就落在由此两区间围成的矩形内. 此矩形显然包含了置信椭圆, 虽然置信椭圆更小些, 但它却以 0.95 的概率包含了均值向量 $\boldsymbol{\mu}=(\mu_{1},\mu_{2})$. 因此, 对由 $T^{2}$ 区间所围成的矩形而言, 其包含两个均值分量 $\mu_{1}$ 与 $\mu_{2}$ 的概率要大于 0.95. 这一结果促使我们去考虑第二种进行多重比较的方法, 即所谓的庞弗罗尼方法.

\mvsub{多重比较的庞弗罗尼方法}

人们常把注意力集中在少数的单一置信区间上, 在这些情形中, 可能会得到比结论 5.3 的联合置信区间更好的结论. 但如果均值分量 $\mu_{i}$ 的个数 $m$ 很小, 或者其线性组合 $\mathbf{a}_{i}'\boldsymbol{\mu}=a_{i1}\mu_{1}+a_{i2}\mu_{2}+\cdots+a_{ip}\mu_{p}$ 很小, 则可采用比联合 $T^{2}$ 区间更为精确的联合置信区间的方法. 另一种多重比较法就是庞弗罗尼方法, 它是由庞弗罗尼概率不等式引申而来的.

假设在得到数据之前, 已知 $m$ 个线性组合 $\mathbf{a}_{1}'\boldsymbol{\mu},\mathbf{a}_{2}'\boldsymbol{\mu},\cdots,\mathbf{a}_{m}'\boldsymbol{\mu}$ 的置信区间, 记 $C_{i}$ 为 $\mathbf{a}_{i}'\boldsymbol{\mu}$ 的置信表示, 满足 $P(C_{i}\text{ 为真})=1-\alpha_{i}, i=1,2,\cdots,m$, 我们有(见练习 5.6)
\[
P(\text{所有 }C_{i}\text{ 为真})=1-P(\text{至少有一个 }C_{i}\text{ 为假})
\]
